\subsection{Экстремальные задачи}

%%%%%%%%%%%%%%%%%%%%%%%%%%%%%%%%%%%%%%%%%%%%%%%%%%%%
%%%%%%%%%%%%%%%%%%%%%%%%%%%%%%%%%%%%%%%%%%%%%%%%%%%%
 
 \z{Необходимо с земли перебросить мяч через вертикальную стенку высоты $h$, находящуюся на расстоянии $s$ от места броска. При какой наименьшей начальной скорости $v_0$ это возможно? Под каким углом $\alpha$ к горизонту должна быть в этом случае направлена начальная скорость мяча?}
%
{$v_0=\sqrt{g(\sqrt{S^2+h^2}+h)};~\alpha=\frac{\pi}{4}+\frac{1}{2}\arctg{\frac{h}{s}}$}

%%%%%%%%%%%%%%%%%%%%%%%%%%%%%%%%%%%%%%%%%%%%%%%%%%%%

 \z{Из миномета, находящегося в точке $A$, ведут обстрел объекта, расположенного на склоне горы. Угол наклона горы к горизонту равен $\beta$, стрельба производится под углом $\alpha$ к горизонту. При каком угле $\alpha_0$ дальность стрельбы вдоль склона будет максимальной? Чему равна эта дальность? }
%
{$\alpha = \frac{\pi}{4}+\frac{\beta}{2};~ L = \frac{2v_0^2\sin{(\alpha-\beta)}\cos{\beta}}{g\cos^2{\beta}}$}
	
%%%%%%%%%%%%%%%%%%%%%%%%%%%%%%%%%%%%%%%%%%%%%%%%%%%%
\z{При какой минимальной начальной скорости $v_\text{мин}$ мальчик может перебросить камень через стену высоты $H$ и толщиной $s$?}
%
{$v_\text{мин}=\sqrt{g(s+2H)}$}

%%%%%%%%%%%%%%%%%%%%%%%%%%%%%%%%%%%%%%%%%%%%%%%%%%%%	
\z[{\begin{wrapfigure}[6]{r}{0.3\textwidth}%
		\vspace{-10pt}%
		\begin{flushright}%
 			\includegraphics[width=0.3\textwidth]{2.5.4}%
		\end{flushright}%
	\end{wrapfigure}%
	}]%
	{Орудие стреляет из-под укрытия, наклоненного к горизонту под углом $\beta$, находясь на расстоянии $L$ основания укрытия. Ствол орудия закреплен под углом $\alpha$ к горизонту, причем $\alpha > \beta$ (см. рисунок). С какой максимальной скоростью $v_\text{макс}$  может вылететь снаряд, не задев укрытия? Сопротивлением воздуха пренебречь. Траектория снаряда лежит в плоскости чертежа.}
%
{$v_\text{макс}=\frac{\sqrt{2gL\tg \beta}}{\cos \alpha (\tg \alpha -\tg \beta)}$}

%%%%%%%%%%%%%%%%%%%%%%%%%%%%%%%%%%%%%%%%%%%%%%%%%%%%	
\z[{\begin{wrapfigure}[6]{r}{0.3\textwidth}%
		\vspace{-15pt}%
		\begin{flushright}%
 			\includegraphics[width=0.3\textwidth]{2.5.5}%
		\end{flushright}%
	\end{wrapfigure}%
	}]%
	{При какой минимальной начальной скорости $v_\text{мин}$ мальчик может перебросить камень через дом с покатой крышей (см. рисунок), если ближайшая к мальчику стена имеет высоту $H$, задняя стена — высоту $h$, а ширина дома равна $L$?}
%
{$v_\text{мин}=\sqrt{g\left(\sqrt{(H-h^2)+L^2}+H+h \right)}$}
